\label{exp:cem00}

\subsubsection*{Постановка}
\textbf{Физический вопрос.}
На микроуровне~$M$ каузальность ограничена тактовой скоростью звена $c_0=\caStepSpace/\htick$;
в учебниках и CODATA фигурирует макроскопическая $\speedC$ и continuum-упаковка $t_P=\ell_P/c$.
Смешиваются ли эти объекты, если задать SI-мост \emph{снизу вверх} от геометрии FCC, а не подгонкой $\kappa=c/c_0$?

\textbf{Что проверяем.}
\begin{enumerate}[label=\textbf{(\alph*)},leftmargin=*,itemsep=0.25em]
\item Одна безразмерная константа $\caKappaGeom$ одновременно связывает шаг времени, шаг энергии и макро-$c$ (\cref{def:si-bridge,sec:fcc-kappa}).
\item Отношение $c_0/c$ следует из геометрии ($\caKappaGeom=1/\sqrt{2}$ на FCC), а не вводится как определение $\kappa$ (\cref{prop:c0-over-c,rem:two-c-burden}).
\item Планковский тик continuum ($t_P$ при $c$) \emph{не} тождественен одному шагу~$g$ на~$M$ (\cref{prop:m-one-tick-cone,def:two-speeds-formal}).
\end{enumerate}

\textbf{Рабочая гипотеза.}
При фиксированных CODATA-константах алгебра SI-моста замыкается: $\kappa_{\mathrm{geom}}=hT/t_P=E_0/E_P$ и $c_0/c=1/\kappa_{\mathrm{geom}}=\sqrt{2}$ с остатками только машинного округления.

\textbf{Наблюдаемые.}
$\kappa_{\mathrm{geom}}$, $hT/t_P$, $c_0/c$, $E_0/E_P$ (табл.~\ref{tab:exp-ce-m-00}) --- четыре проекции одной строки констант.

\textbf{Критерий согласия с теорией.}
Совпадение трёх проекций $\kappa$; $|c_0/c-\sqrt{2}|<10^{-6}$; алгебраические остатки SI-тождеств $<10^{-12}$ (\cref{sec:compare-reference}, тип~B).

\subsubsection*{Теоретическая часть}
Один такт~$g$ на~$M$ --- один шаг по $N_{12}$ внутри светового конуса (\cref{prop:m-one-tick-cone}).
Тактовая скорость $c_0$ и скорость, с которой макроскопический фронт после осреднения сопоставляется с $\speedC$, --- разные уровни описания (\cref{def:two-speeds-formal}).
SI-мост \emph{определяет} $\htick=\caKappaGeom t_P$ и $\speedC=\caKappaGeom c_0$, где $\caKappaGeom$ выводится из изотропии FCC (\cref{sec:fcc-kappa}), а не из измеренного отношения $c/c_0$.

Опыт чисто алгебраический: поле $z$ на решётке не эволюционирует; проверяется согласованность определений перед любыми гидро- и SM-сверками.

\subsubsection*{Ход эксперимента}
Модуль констант \texttt{mt\_ca.si\_constants.SI}; воспроизведение:
\texttt{python verify\_principles.py --suite carrier}.

\begin{enumerate}[label=\textbf{Шаг \arabic*:},leftmargin=*,itemsep=0.35em]
\item Зафиксировать коммит и окружение (\cref{ch:computational-laboratory}).
\item Вычислить строку SI-моста и перенести значения в табл.~\ref{tab:exp-ce-m-00}.
\item Сравнить $c_0/c$ с $1/\kappa_{\mathrm{geom}}$ и с $\sqrt{2}$.
\end{enumerate}

\begin{table}[htbp]
\centering
\small
\caption{Наблюдаемые SI-моста (CE-M-00)}
\label{tab:exp-ce-m-00}
\begin{tabular}{@{}ll@{}}
\toprule
Величина & Значение \\
\midrule
$\kappa_{\mathrm{geom}}$ (FCC) & $0{,}707107$ \\
$hT/t_P$ & $0{,}707107$ \\
$c_0/c$ & $1{,}414214$ \\
$E_0/E_P$ & $0{,}707107$ \\
\bottomrule
\end{tabular}
\end{table}

\subsubsection*{Анализ, расчёт погрешности и выводы}
Три проекции $\kappa$ совпадают побитово в таблице: геометрия FCC, время и энергия на звене --- одна константа.
$c_0/c=\sqrt{2}$ --- следствие $\caKappaGeom$, не подгонка (\cref{prop:c0-over-c}).

Погрешность: тип~B (алгебра + float); ансамбль не требуется.

\textbf{Вывод.}
Continuum-формула $t_P=\ell_P/c$ описывает \emph{макро}упаковку; микро-тик и $c_0$ живут на~$M$ и связываются с CODATA только через явный мост.
\textbf{Статус записи:} опыт закрыт по согласованию с \cref{rem:two-c-burden}.
